% sections/08_related.tex --- SkillStrike

\paragraph{Prompt injection.}
Greshake \etal~\cite{greshake2023indirect} introduce \emph{indirect prompt injection},
where attacker-controlled web content manipulates an LLM agent's behavior.
Perez and Ribeiro~\cite{perez2022ignore} study direct prompt injection in single-turn settings.
Our work differs: we do not inject adversarial inputs into the LLM directly;
we exploit the agent's \emph{trusted tool execution} path by poisoning
the output of a co-installed source skill.
The agent believes it is acting on legitimate skill output,
not adversarial user input.
This distinction means prompt-injection defenses (input sanitization,
constrained decoding) are structurally inapplicable to our attack.

\paragraph{LLM agent security.}
Ruan \etal~\cite{ruan2024identifying} characterize ambient authority risks
in LLM agent systems and propose a safety evaluation framework.
Zhan \etal~\cite{zhan2024injecagent} study tool-call injection in multi-agent pipelines.
Yi \etal~\cite{yi2023benchmarking} benchmark LLM safety against prompt attacks.
Our contribution is orthogonal: we study \emph{compositional risk} arising from
the combination of skills, not the behavior of individual skills or models in isolation.

\paragraph{Software supply-chain attacks.}
Zimmermann \etal~\cite{zimmermann2019small} and Synk~\cite{synk2023report}
study malicious package injection in npm and PyPI.
Xi \etal~\cite{xi2023rise} analyze tool-related risks in LLM ecosystems broadly.
Prior supply-chain attacks target code execution at install or import time.
Our attacks execute \emph{at agent invocation time},
require no malicious code in the skill itself,
and leverage the agent's own tool-call machinery rather than OS-level execution.

\paragraph{Compositional security in traditional systems.}
The principle that composition increases attack surface
is well-studied in traditional access control~\cite{felt2011android},
microservice security, and inter-program data flow~\cite{enck2014taintdroid}.
Our work brings this principle to the new domain of LLM agent ecosystems,
where the composition is not statically declared, not sandboxed, and
not visible to the user at install time.

\paragraph{MCP and agent-tool security.}
Concurrent work analyzes MCP server trust~\cite{anthropic2024mcp},
tool schema injection, and server identity spoofing.
Our attack uses identity spoofing (\S\ref{sec:eval_ch}) as one mechanism,
but the core vulnerability is the \emph{shared context without capability isolation},
which exists regardless of whether source skill output is spoofed or legitimate.
